\documentclass[11pt]{article}
\usepackage[a4paper,margin=1in]{geometry}
\usepackage{amsmath,amssymb,amsthm}
\usepackage[T1]{fontenc}
\usepackage{lmodern}
\usepackage[hidelinks]{hyperref}
\setlength{\emergencystretch}{3em}

\theoremstyle{plain}
\newtheorem{theorem}{Theorem}
\newtheorem{lemma}[theorem]{Lemma}
\theoremstyle{definition}
\newtheorem{definition}[theorem]{Definition}
\theoremstyle{remark}
\newtheorem{remark}[theorem]{Remark}

\title{A Counterexample to Conjecture 00000001666:\\
the Game Domination Number of $C_n$ Is at Least $n/3$}
\author{%
  Feiyu\thanks{Independent researcher.}
  \and
  Claude (Opus 5.5)\thanks{Anthropic.}%
}
\date{2026-10-02}

\begin{document}
\maketitle

\begin{abstract}
Conjecture 00000001666 asserts that the game domination number of the cycle
$C_n$ is $\lceil n/5\rceil+O(1)$. Whatever the players do, the domination game on
$C_n$ lasts at least $n/3$ moves: a vertex of $C_n$ dominates at most three
vertices, and all $n$ vertices must be dominated. Since $n/3-n/5$ is unbounded,
the conjecture is false. The game, its value under optimal play and the bound
are formalized in Lean~4 against Mathlib.
\end{abstract}

\section{The conjecture as stated}

The original text of conjecture 00000001666 reads:

\begin{quote}
Definition: The game domination number $\gamma_g(G)$ is the two-player
alternating-choice game version. Conjecture: The game domination number of the
cycle $C_n$ is $\lceil n/5\rceil + O(1)$, with the constant term explicit by $n
\bmod 5$ (cycle game domination).
\end{quote}

\section{Reading of the statement}

We use the standard domination game of Bre\v{s}ar, Klav\v{z}ar and Rall
\cite{bkr}.

\begin{definition}
In the \emph{domination game} on a finite graph $G$, two players, Dominator and
Staller, alternately choose vertices, Dominator first. Each chosen vertex must
dominate (be equal or adjacent to) at least one vertex not dominated by the
previously chosen ones. The game ends when every vertex is dominated. Dominator
aims to minimize the number of moves, Staller to maximize it, and the
\emph{game domination number} $\gamma_g(G)$ is the number of moves under optimal
play.
\end{definition}

The conjecture asserts that there is a constant $C$ with
$|\gamma_g(C_n)-\lceil n/5\rceil|\le C$ for all $n\ge3$. (The clause on the
constant term depending on $n\bmod5$ only refines this.)

\section{Result}

\begin{theorem}\label{thm:main}
For every $n$, $\gamma_g(C_n)\ge n/3$. Consequently
$\gamma_g(C_n)-\lceil n/5\rceil\ge n/3-n/5-1$ is unbounded, and Conjecture
00000001666 is false.
\end{theorem}

\section{Proof}

\begin{lemma}\label{lem:game}
Let $G$ be a finite graph in which every closed neighbourhood $N[v]$ has at most
$m$ vertices. In any position of the domination game, under any play, at least
$u/m$ further moves are made, where $u$ is the number of undominated
vertices. In particular $\gamma_g(G)\ge |V(G)|/m$.
\end{lemma}

\begin{proof}
Each move dominates only vertices of $N[v]$ for the chosen $v$, hence lowers the
number of undominated vertices by at most $m$. The game ends only when this
number is $0$. Formally, by induction on the game tree: if a legal move exists,
the value is $1$ plus the minimum (Dominator) or maximum (Staller) of the values
after the legal moves, each of which is at least $(u-m)/m$ by induction. If no
legal move exists, all vertices are dominated and $u=0$.
\end{proof}

\begin{proof}[Proof of Theorem~\ref{thm:main}]
In $C_n$ every closed neighbourhood has at most $3$ vertices, so
Lemma~\ref{lem:game} with $m=3$ gives $\gamma_g(C_n)\ge n/3$. If
$\gamma_g(C_n)\le\lceil n/5\rceil+C$ for all $n\ge3$, take $n=15C+15$. Then
$5C+5\le\gamma_g(C_n)\le 3C+3+C$, which is impossible.
\end{proof}

\section{Formalization}

The accompanying Lean~4 project (\texttt{lean/Submission/Basic.lean}) defines
the game, states the conjecture and proves its negation against Mathlib.

\begin{itemize}
  \item \texttt{closedNbhd}, \texttt{dominated}, \texttt{legalMoves} --- for
    any finite \texttt{SimpleGraph} with decidable adjacency.
  \item \texttt{gameValue fuel dom S} --- the value of the game from position $S$
    (chosen vertices), $0$ when no legal move remains, otherwise $1+$ the
    minimum (Dominator, \texttt{dom = true}) or maximum (Staller) over legal
    moves. The parameter \texttt{fuel} bounds the remaining moves; since every
    move dominates a new vertex, \texttt{fuel} $=|V|$ is never exhausted.
    \texttt{gameDominationNumber G} $=$ \texttt{gameValue G |V| true} $\emptyset$.
  \item \texttt{card\_undominated\_le\_mul},
    \texttt{card\_le\_mul\_gameDominationNumber} --- Lemma~\ref{lem:game}.
  \item \texttt{card\_closedNbhd\_cycleGraph},
    \texttt{le\_three\_mul\_gameDominationNumber} --- $n\le3\gamma_g(C_n)$ for
    Mathlib's \texttt{cycleGraph n}.
  \item \texttt{ConjectureHolds} --- there is $C$ with
    $|\gamma_g(C_n)-\lceil n/5\rceil|\le C$ for all $n\ge3$;
    \texttt{conjecture\_00000001666\_false} --- its negation.
\end{itemize}

The toolchain is \texttt{leanprover/lean4:v4.33.1}, Mathlib is pinned to
\texttt{v4.33.1}, and \texttt{lake build} succeeds. The project contains no
\texttt{sorry}, no \texttt{admit} and no \texttt{native\_decide}, and
introduces no axioms:
\begin{center}
\texttt{\#print axioms conjecture\_00000001666\_false}
\end{center}
reports exactly \texttt{propext}, \texttt{Classical.choice} and
\texttt{Quot.sound}.

\section{Remarks}

The game domination number of cycles is in fact about $n/2$. It was determined
exactly by Ko\v{s}mrlj \cite{kosmrlj}. The elementary bound $n/3$, which follows
from $\gamma_g(G)\ge\gamma(G)$, already refutes the conjectured growth
$\lceil n/5\rceil+O(1)$.

\begin{thebibliography}{9}
\bibitem{bkr} B.~Bre\v{s}ar, S.~Klav\v{z}ar and D.~F.~Rall, \emph{Domination
  game and an imagination strategy}, SIAM J. Discrete Math. 24 (2010),
  979--991.
\bibitem{kosmrlj} G.~Ko\v{s}mrlj, \emph{Domination game on paths and cycles},
  Ars Math. Contemp. 13 (2017).
\end{thebibliography}

\end{document}
